\documentclass[11pt,a4paper]{article}
\usepackage[T1]{fontenc}
\usepackage[utf8]{inputenc}
\usepackage[polish]{babel}
\usepackage{lmodern}
\usepackage{microtype}
\usepackage[a4paper,margin=2.25cm]{geometry}
\usepackage{amsmath,amssymb,mathtools,bm}
\usepackage{booktabs}
\usepackage{hyperref}
\title{Dodatek matematyczny\\Model warunkowany horyzontem dla generycznych parametrów}
\author{Edward Szczypka}
\date{\today}
\newcommand{\E}{\mathbb{E}}
\newcommand{\R}{\mathbb{R}}
\begin{document}
\maketitle

\section{Rodzina parametrów}

Niech $\mathcal P$ będzie skończonym, konfigurowalnym zbiorem parametrów
jakości powietrza, np.
\[
 \mathcal P=\{\mathrm{PM10},\mathrm{PM2.5},\mathrm{NO2},\mathrm{O3}\}.
\]
Dla $p\in\mathcal P$, stacji $s$, czasu bazowego $t$ i horyzontu
$h\in\{1,\ldots,H\}$ zmienna docelowa ma postać:
\[
 Y^{(p)}_{s,t,h}=X^{(p)}_{s,t+h}.
\]

\section{Model wspólnego kontraktu}

Każdy parametr może mieć inny provider, lecz wspólny kontrakt:
\[
 \widehat Y^{(p)}_{s,t,h}
 =f_{p,\theta_p}\!\left(
      \bm x^{(p)}_{s,t},
      h,
      \bm\phi(t+h),
      \bm w_{s,t+h},
      \bm z^{(p)}_{s,t}
   \right),
\]
gdzie $\bm w$ oznacza prognozowane cechy meteorologiczne, a
$\bm z^{(p)}$ cechy innych parametrów skonfigurowanych jako pomocnicze.

\section{Cechy pomocnicze i kauzalność}

Dla $q\in\mathcal A_p\subseteq\mathcal P\setminus\{p\}$ stosujemy:
\[
 \bm z^{(q)}_{s,t}
 =\left(X^{(q)}_{s,t},X^{(q)}_{s,t-1},X^{(q)}_{s,t-6},X^{(q)}_{s,t-24}\right).
\]
Warunek kauzalności:
\[
 \sigma(\bm x_{s,t},\bm z_{s,t})\subseteq\mathcal F_t,
\]
czyli wszystkie cechy muszą być dostępne w chwili tworzenia prognozy.

\section{Ograniczenia wartości}

Definicja parametru zawiera zbiór dopuszczalny
\[
 D_p=[a_p,b_p],
\]
gdzie koniec może być nieskończony. Publikowana predykcja jest projekcją:
\[
 \widetilde Y^{(p)}=\Pi_{D_p}(\widehat Y^{(p)}).
\]
Zapobiega to wspólnemu, błędnemu założeniu, że wszystkie parametry mają ten sam
zakres lub jednostkę.

\section{Ważona funkcja celu}

Dla próbki $\mathcal I_p$:
\[
 \mathcal L_p(\theta_p)
 =\sum_{i\in\mathcal I_p}w_i\,
   \ell_p\!\left(y_i,f_{p,\theta_p}(x_i,h_i)\right),
\]
przy czym waga może uwzględniać świeżość, stację, horyzont i rzadkie epizody.
Provider wybiera $\ell_p$ zgodnie z charakterem parametru; rejestr nie wymusza
jednego estymatora ani jednej funkcji straty.

\section{Brama jakości}

Kandydat parametru $p$ zastępuje model aktywny wyłącznie, gdy na późniejszym
okresie chronologicznym:
\[
 \frac{\operatorname{MAE}_{\mathrm{baseline},p}
       -\operatorname{MAE}_{\mathrm{candidate},p}}
      {\operatorname{MAE}_{\mathrm{baseline},p}}
 \ge \eta_p.
\]
Wartość $\eta_p$ może być parametryzowana oddzielnie dla każdego parametru.

\section{Skalowalność}

Dodanie parametru do katalogu nie zwiększa kosztu ML. Koszt powstaje dopiero po
włączeniu roli celu. Dla zbioru celów $\mathcal T$ koszt treningu jest w
przybliżeniu:
\[
 C\approx\sum_{p\in\mathcal T}
   N_p\,K_p\,F_p,
\]
gdzie $N_p$ jest ograniczoną liczbą rekordów, $K_p$ liczbą kandydatów, a
$F_p$ liczbą foldów. Profile quick/full ograniczają każdy z tych czynników.

\end{document}
